% =====================================================================
% EgoAnchor — IEEE VR 2027 最终版主稿（深度润色版）
% 基于 egoanchor_final.tex 重构，摘要精简至150词，引言深度润色
% 编译产物输出到 2026-EgoAnchor/pdf/
% 参考文献：egoanchor_cn_refs.bib
% =====================================================================
\documentclass[review]{vgtc}

\graphicspath{{figures/}{pictures/}{images/}{./}}

\usepackage[UTF8]{ctex}
\usepackage{times}
\renewcommand*\ttdefault{txtt}
\usepackage{mathptmx}
\usepackage{booktabs}
\usepackage{tabularx}
\usepackage{multirow}
\usepackage{enumitem}
\usepackage{amsmath}
\usepackage{amssymb}
\usepackage{graphicx}
\usepackage{xcolor}

\onlineid{0}
\vgtccategory{Research}
\vgtcinsertpkg

\title{EgoAnchor: A Training-Free Real-Object Anchoring System \\ for Passthrough Mixed Reality on Consumer Hardware}

\author{Anonymous for Review}
\authorfooter{
% 匿名审稿阶段；正式投稿时填入作者信息
}

\teaser{
  \centering
  \fbox{\rule{0pt}{2.0in}\rule{0.95\linewidth}{0pt}}
  \caption{EgoAnchor system overview and capability positioning. \textbf{(Top)} Capability matrix: no existing system simultaneously satisfies all five dimensions (depth-free, arbitrary objects, dynamic tracking, training-free, open hardware). \textbf{(Middle)} End-to-end architecture: Quest~3 binocular capture → external tracking service (segmentation/depth/6DoF estimation) → four-layer coordinated runtime (temporal alignment / quality gating / smoothing / lifecycle management). \textbf{(Bottom)} Comparison: traditional arrival-time mapping causes slip; EgoAnchor achieves stable attachment. Placeholder: suggest final figure showing (a) capability comparison table, (b) system architecture diagram, (c) slip vs. stable attachment comparison.}
  \label{fig:teaser}
}

\abstract{%
Stable anchoring of virtual content to physical objects is fundamental to passthrough mixed reality, yet existing solutions face a capability coverage gap. Platform-specific approaches either require proprietary depth sensors (HoloLens, Vision~Pro), limit to static objects (Vuforia Model Targets), or support only predefined categories (Meta Dynamic Object Tracker). We present \textbf{EgoAnchor}, the first system simultaneously achieving: training-free deployment, binocular-only sensing, arbitrary rigid daily objects, dynamic tracking, and open consumer hardware. Our key insight is that world-consistent anchoring demands systematic coordination across time alignment, quality gating, temporal smoothing, and lifecycle management---not merely accurate per-frame pose estimation. EgoAnchor introduces a four-layer architecture: (1)~frame-aligned anchoring decouples perception latency from head motion by retrieving capture-time camera poses; (2)~reliability gating filters anomalous observations via multi-dimensional scoring; (3)~Kalman-based temporal smoothing with static lock suppresses jitter; (4)~lifecycle management enables rapid reacquisition after occlusion. Operating at 5--12\,fps perception input, the system produces 60\,fps world-consistent anchors on consumer GPUs. Evaluation on Quest~3 demonstrates anchoring quality comparable to platform solutions, establishing a viable open pathway for MR physical interaction. Code will be released.}

\keywords{Mixed Reality, Real-Object Anchoring, 6DoF Tracking, Asynchronous Perception, System Design, Consumer Hardware}

\begin{document}

\maketitle

% =====================================================================
% 1. 引言（深度润色，达到顶会水准）
% =====================================================================
\section{Introduction}

Passthrough mixed reality promises to seamlessly blend virtual content with the physical world, enabling applications from spatially-anchored annotations to hand-held object augmentation. A fundamental requirement for such interactions is \emph{stable anchoring}: virtual overlays must remain perceptually locked to their associated physical objects as users move their heads, approach targets, and manipulate devices. Achieving this stability requires systems to satisfy multiple capability dimensions simultaneously, yet existing solutions make inherent trade-offs that leave critical gaps in coverage.

\begin{table}[t]
\centering
\caption{Capability coverage of real-object anchoring approaches. EgoAnchor is the first to satisfy all five dimensions simultaneously.}
\label{tab:capability-comparison}
\small
\begin{tabular}{@{}lccccc@{}}
\toprule
\textbf{Approach} & \makecell{\textbf{Depth-}\\\textbf{Free}} & \makecell{\textbf{Arbitrary}\\\textbf{Objects}} & \makecell{\textbf{Dynamic}\\\textbf{Tracking}} & \makecell{\textbf{Training-}\\\textbf{Free}} & \makecell{\textbf{Open}\\\textbf{Hardware}} \\
\midrule
Azure Object Anchors~\cite{azureObjectAnchors} & \xmark & \cmark & \cmark & \xmark & \xmark \\
Vision Pro Object Tracking~\cite{appleVisionOS27ObjectTracking} & \xmark & \cmark & \cmark & \xmark & \xmark \\
Vuforia Model Targets~\cite{vuforiaModelTargets} & \cmark & \cmark & \xmark & \cmark & \xmark \\
Meta Dynamic Object Tracker~\cite{metaDynamicObjectTracker} & \cmark & \xmark & \cmark & \cmark & \xmark \\
\midrule
\textbf{EgoAnchor (ours)} & \cmark & \cmark & \cmark & \cmark & \cmark \\
\bottomrule
\end{tabular}
\end{table}

Active-depth solutions provide high-quality 3D sensing but require proprietary sensors. Azure Object Anchors on HoloLens~2 leveraged spatial mapping but was discontinued in late 2023~\cite{azureObjectAnchors}. Apple Vision Pro extended object tracking to hand-held and moving objects in visionOS~27~\cite{appleVisionOS27ObjectTracking}---a significant advance---yet it remains tied to LiDAR-equipped closed hardware, requires hours of Create~ML offline training per object, and reportedly achieves only ~30\,Hz even on its high-rate pathway. Camera-only solutions such as Vuforia Model Targets achieve markerless tracking via edge matching but are explicitly scoped to \emph{static} objects; movement breaks registration~\cite{vuforiaModelTargets}. When platforms do support moving objects, coverage contracts to a handful of predefined categories: Meta Quest's dynamic object tracker stably handles only keyboards and similar fixed types, with documentation noting that update rates are unsuitable for real-time following of moving objects~\cite{metaDynamicObjectTracker,metaSceneOverview}. The combination of \textbf{depth-free sensing, arbitrary daily objects, dynamic tracking, training-free deployment, and open consumer hardware} remains unaddressed (Table~\ref{tab:capability-comparison}).

We present \textbf{EgoAnchor}, an end-to-end system that closes this five-dimensional gap. Operating solely from binocular passthrough images and object 3D models---obtainable from casual phone captures via image-to-3D generation~\cite{trellis2024,appleObjectCapture}---EgoAnchor anchors arbitrary rigid daily objects (mice, headphones, controllers) on consumer GPUs without per-object training. Our central thesis is that world-consistent anchoring is \emph{not} guaranteed by per-frame camera-space pose accuracy alone; it demands systematic coordination across time alignment, quality awareness, temporal smoothing, and lifecycle recovery. The challenge is to transform a low-rate (5--12\,fps), noisy, intermittently-failing asynchronous perception stream into high-rate (60\,fps), world-consistent, stable anchors suitable for interactive MR applications.

The core challenge is \textbf{temporal asynchrony}. To achieve reliable 6DoF estimation, systems offload neural network inference to external compute; inference latency renders perception low-rate and asynchronous, while HMD rendering must sustain $>$90\,fps to avoid discomfort. Naively multiplying a returned camera-space pose by the \emph{arrival-time} latest HMD pose to map into world space couples any head motion during the capture-to-return interval into the result, manifesting as visible \emph{slip} of virtual content relative to the physical object. This phenomenon echoes the dynamic registration error long identified in AR systems---Azuma noted that latency multiplied by motion is the primary contributor~\cite{azuma1997survey}---but occurs in the \emph{perception fusion} stage rather than the render-to-display pathway. Beyond temporal misalignment, vision-based tracking inevitably produces noise and pose jumps under motion blur, lighting changes, or partial occlusions; directly driving virtual content with unfiltered observations causes severe visual jitter.

EgoAnchor addresses these challenges through a \textbf{four-layer coordinated architecture}:

\noindent\textbf{Layer 1: Temporal Alignment.} \emph{Frame-aligned anchoring} binds each outgoing image frame with a unique \texttt{frame\_id} and the capture-time camera world pose cached in a ring buffer. When the external service returns a 6DoF observation for that frame, the Unity frontend retrieves the \emph{capture-time} cached pose rather than using the current latest HMD pose, ensuring world-space mapping corresponds to the physically correct instant. This design mirrors the principle of late latching / asynchronous timewarp in rendering pipelines~\cite{vanwaveren2016atw,mark1997postrendering}---timewarp aligns rendered output to the \emph{current} latest pose (compensating render latency), while frame-aligned anchoring aligns perception output back to its \emph{capture instant} (compensating perception latency). Both align temporally-tagged data to the physically correct moment rather than blindly multiplying across latency, but frame-aligned anchoring uses \emph{recorded history} rather than predicted futures, eliminating slip via precision rather than estimation.

\noindent\textbf{Layer 2: Quality Gating.} Multi-dimensional reliability scoring---combining LAB-space color reprojection error and rendered depth validity within masks---gates observations, blocking anomalous poses from polluting downstream stages.

\noindent\textbf{Layer 3: Temporal Smoothing.} A constant-velocity Kalman filter adaptively estimates object world-space position, velocity, and rotation. When relative motion falls below thresholds for a settling period, a \emph{static lock} controller engages, clamping output to suppress residual micro-jitter via dead-zone CUSUM and drift leash mechanisms, with seamless unlocking and seam removal upon motion resumption.

\noindent\textbf{Layer 4: Lifecycle Management.} A state machine (Searching / Tracking / Coasting / Lost) manages occlusion and failure: Coasting maintains bounded Kalman prediction during transient occlusion with visual feedback, Lost triggers backend re-registration requests via a command channel, achieving rapid self-healing recovery.

We emphasize that individual layer components do not claim algorithmic novelty; their contribution lies in being systematically organized into a coordinated runtime tailored for asynchronous MR perception. We implemented and evaluated EgoAnchor on Quest~3 with consumer GPUs (RTX~5080 Laptop / 5090 Desktop), establishing an anchor-centric evaluation protocol that shifts focus from camera-space pose accuracy to application-layer anchor quality: world-space error, static jitter, head-motion-induced slip, end-to-end latency, and recovery time. Results demonstrate anchoring quality comparable to platform-level solutions, providing a viable open pathway for cross-platform MR physical interaction systems.

\noindent\textbf{Contributions:}
\begin{enumerate}[leftmargin=1.4em]
  \item \textbf{First five-dimensional system.} EgoAnchor simultaneously achieves depth-free sensing, arbitrary objects, dynamic tracking, training-free deployment, and open hardware---a combination no existing system provides.
  \item \textbf{Four-layer coordinated architecture.} We design and implement a layered runtime (temporal alignment, quality gating, smoothing, lifecycle) that transforms asynchronous perception streams into world-consistent, stable, recoverable anchors.
  \item \textbf{Anchor-centric evaluation protocol.} We establish metrics centered on anchor quality (world error, jitter, slip, latency, recovery) rather than per-frame pose accuracy, validated on Quest~3 with consumer hardware.
  \item \textbf{Open implementation.} Complete system code will be released, demonstrating that open systems can achieve platform-comparable anchoring quality on consumer hardware.
\end{enumerate}

% =====================================================================
% 后续章节：提纲形式，供后续填充
% =====================================================================

\section{Related Work}

\subsection{Platform-Level Real-Object Anchoring}
% 提纲：对比五维能力空间，展开表格细节，说明为何现有方案各有取舍

\subsection{Markerless 6DoF Object Pose Estimation and Tracking}
% 提纲：CV领域的 register/track/re-register 范式，FoundationPose 等基础模型，评价以 ADD/ADD-S 为中心，不涉及世界一致性

\subsection{XR Latency, Registration, and Temporal Filtering}
% 提纲：AR 配准误差源头（Azuma 1997），预测式跟踪（Azuma & Bishop 1994），时间扭曲（van Waveren 2016），One Euro / Kalman 滤波折衷

\section{System Design}

\subsection{Design Goals and System Challenges}
% 提纲：世界一致性、稳定可用性、可恢复性、可分析性四个目标；时间异步挑战的数学表述

\subsection{Architecture Overview}
% 提纲：前后端分离架构，Quest/Unity 采集 → ZMQ 数据面 → Python 追踪 → NATS 消息/命令面 → Unity 四层运行时

\section{Four-Layer Coordinated Anchoring Architecture}

\subsection{Object-Level Perception Frontend}
% 提纲：相机参数映射、prompt 分割（YOLOE/SAM3）、双目深度（FoundationStereo）、register/track/re-register（FoundationPose）

\subsection{Layer 1: Temporal Alignment---Frame-Aligned Anchoring}
% 提纲：frame_id 回查机制，与时间扭曲的对偶关系，坐标系对齐

\subsection{Layer 2: Quality Gating---Reliability Scoring}
% 提纲：多维评分（颜色重投影 + 深度有效性），几何对数平均公式，门控逻辑

\subsection{Layer 3: Temporal Smoothing---Kalman Filter and Static Lock}
% 提纲：恒速卡尔曼状态估计，静止锁定（进入条件、三路解锁、接缝消除）

\subsection{Layer 4: Lifecycle Management---Hold and Reacquisition}
% 提纲：Searching/Tracking/Coasting/Lost 状态机，NATS 命令面重获取

\section{Implementation}

\subsection{Hardware and Software Environment}
% 提纲：Quest~3、RTX~5080/5090、Unity/Python/CUDA 版本

\subsection{Network and Communication}
% 提纲：ZMQ Protobuf 数据面、NATS Protobuf 消息面与命令面

\subsection{Offline Simulation and Diagnostics}
% 提纲：Tools3 离线回放、参数敏度分析

\section{Evaluation Design}

\subsection{Research Questions}
\begin{itemize}[leftmargin=1.4em]
  \item \textbf{RQ1 (System Performance):} Can EgoAnchor transform asynchronous perception streams into world-consistent, stable anchors on consumer hardware?
  \item \textbf{RQ2 (Architecture Contribution):} What are the ablation benefits of the four-layer architecture in eliminating slip, suppressing jitter, and controlling latency?
  \item \textbf{RQ3 (Robustness \& Recovery):} How do recovery success rate and recovery time perform under occlusion, out-of-view, and fast head motion?
  \item \textbf{RQ4 (Task Usability):} Does stable anchoring improve user task performance and subjective trust?
\end{itemize}

\subsection{Baselines and Ablations}
% 提纲：系统级基线（Arrival-Time / Frame-Aligned Raw / EgoAnchor Default）；模块消融（滤波器对比 Pareto 前沿、静止锁定、质量门控）

\subsection{Metrics}
% 提纲：主指标（世界锚定误差、静态抖动、头动打滑）；辅助指标（时延、恢复时间/成功率）；支持性指标（相机坐标系位姿误差）

\subsection{Ground Truth and Experimental Conditions}
% 提纲：Quest Touch 手柄作为主真值、ArUco/AprilTag 用于日常物体；须披露的方法学局限（共用跟踪系、标记翻转歧义、跨设备对比仅定性）

\subsection{User Study Design}
% 提纲：被试内 N=12-15、功效分析、两任务（精细对齐/遮挡恢复）、Likert 量表与统计方法

\section{Results}
% 提纲：对应 RQ1-4 的实验结果，待填充数据与图表

\section{Discussion}

\subsection{System Contribution: From Pose Accuracy to Anchor Quality}
% 提纲：为何单帧精度不足以保证锚点质量，系统协同的本质

\subsection{Five-Dimensional Gap: Ecosystem Implications}
% 提纲：开放系统可达到平台级质量的生态意义

\subsection{Four-Layer Design Principles}
% 提纲：分而治之、各司其职、层间解耦

\subsection{Applicability and Failure Modes}
% 提纲：适用边界（物体特性、运动速度、视野外移动、外部算力依赖）

\subsection{Limitations and Future Work}
% 提纲：多物体扩展、端侧轻量化、学习式时序策略、长时稳定性

\section{Conclusion}
% 提纲：回到五维能力空白，总结四层协同设计如何将异步感知流转化为可用锚点，强调对开放 MR 系统的参考意义

\nocite{*}
\bibliographystyle{abbrv-doi-hyperref-narrow}
\bibliography{egoanchor_cn_refs}

\end{document}
